\documentclass[10pt,a4paper]{amsart}
\usepackage{amsmath,amssymb,amsthm,graphicx}
\usepackage{verbatim}
\newenvironment{reference}{\comment}{\endcomment}
%\newenvironment{reference}{}{}
\newenvironment{slogan}{\comment}{\endcomment}
\newenvironment{history}{\comment}{\endcomment}

% For commutative diagrams we use Xy-pic
\usepackage[all]{xy}

% We use 2cell for 2-commutative diagrams.
\xyoption{2cell}
% Theorem environments.
%
\theoremstyle{plain}
\newtheorem{theorem}[subsection]{Theorem}
\newtheorem{proposition}[subsection]{Proposition}
\newtheorem{lemma}[subsection]{Lemma}

\theoremstyle{definition}
\newtheorem{definition}[subsection]{Definition}
\newtheorem{example}[subsection]{Example}
\newtheorem{exercise}[subsection]{Exercise}
\newtheorem{situation}[subsection]{Situation}

\theoremstyle{remark}
\newtheorem{remark}[subsection]{Remark}
\newtheorem{remarks}[subsection]{Remarks}

\numberwithin{equation}{subsection}

% Macros
%
\def\lim{\mathop{\mathrm{lim}}\nolimits}
\def\colim{\mathop{\mathrm{colim}}\nolimits}
\def\Spec{\mathop{\mathrm{Spec}}}
\def\Hom{\mathop{\mathrm{Hom}}\nolimits}
\def\Ext{\mathop{\mathrm{Ext}}\nolimits}
\def\SheafHom{\mathop{\mathcal{H}\!\mathit{om}}\nolimits}
\def\SheafExt{\mathop{\mathcal{E}\!\mathit{xt}}\nolimits}
\def\Sch{\mathit{Sch}}
\def\Mor{\mathop{\mathrm{Mor}}\nolimits}
\def\Ob{\mathop{\mathrm{Ob}}\nolimits}
\def\Sh{\mathop{\mathit{Sh}}\nolimits}
\def\NL{\mathop{N\!L}\nolimits}
\def\CH{\mathop{\mathrm{CH}}\nolimits}
\def\proetale{{pro\text{-}\acute{e}tale}}
\def\etale{{\acute{e}tale}}
\def\QCoh{\mathit{QCoh}}
\def\Ker{\mathop{\mathrm{Ker}}}
\def\Im{\mathop{\mathrm{Im}}}
\def\Coker{\mathop{\mathrm{Coker}}}
\def\Coim{\mathop{\mathrm{Coim}}}

% Boxtimes
%
\DeclareMathSymbol{\boxtimes}{\mathbin}{AMSa}{"02}

%
% Macros for moduli stacks/spaces
%
\def\QCohstack{\mathcal{QC}\!\mathit{oh}}
\def\Cohstack{\mathcal{C}\!\mathit{oh}}
\def\Spacesstack{\mathcal{S}\!\mathit{paces}}
\def\Quotfunctor{\mathrm{Quot}}
\def\Hilbfunctor{\mathrm{Hilb}}
\def\Curvesstack{\mathcal{C}\!\mathit{urves}}
\def\Polarizedstack{\mathcal{P}\!\mathit{olarized}}
\def\Complexesstack{\mathcal{C}\!\mathit{omplexes}}
% \Pic is the operator that assigns to X its picard group, usage \Pic(X)
% \Picardstack_{X/B} denotes the Picard stack of X over B
% \Picardfunctor_{X/B} denotes the Picard functor of X over B
\def\Pic{\mathop{\mathrm{Pic}}\nolimits}
\def\Picardstack{\mathcal{P}\!\mathit{ic}}
\def\Picardfunctor{\mathrm{Pic}}
\def\Deformationcategory{\mathcal{D}\!\mathit{ef}}
\usepackage{cleveref}
\crefname{originalreference}{original reference}{original references}
\Crefname{originalreference}{Original reference}{Original references}
\makeatletter
\newlabel{lemma-extend-over-normal-curve}{{0BXZ}{source}{}{sourcepointer-1}{}}
\newlabel{lemma-extend-over-normal-curve@cref}{{[originalreference][0][]0BXZ}{[1][0][]source}{}{}{}}
\makeatother
\crefname{originalreference}{original reference}{original references}
\Crefname{originalreference}{Original reference}{Original references}
\makeatletter
\newlabel{lemma-generically-etale}{{0C1C}{source}{}{sourcepointer-2}{}}
\newlabel{lemma-generically-etale@cref}{{[originalreference][0][]0C1C}{[1][0][]source}{}{}{}}
\makeatother
\crefname{originalreference}{original reference}{original references}
\Crefname{originalreference}{Original reference}{Original references}
\makeatletter
\newlabel{lemma-linear-series}{{0CCP}{source}{}{sourcepointer-3}{}}
\newlabel{lemma-linear-series@cref}{{[originalreference][0][]0CCP}{[1][0][]source}{}{}{}}
\makeatother
\crefname{originalreference}{original reference}{original references}
\Crefname{originalreference}{Original reference}{Original references}
\makeatletter
\newlabel{proposition-unwind-morphism-smooth}{{0CD2}{source}{}{sourcepointer-4}{}}
\newlabel{proposition-unwind-morphism-smooth@cref}{{[originalreference][0][]0CD2}{[1][0][]source}{}{}{}}
\makeatother
\crefname{originallemma}{original lemma}{original lemmas}
\Crefname{originallemma}{Original lemma}{Original lemmas}
\makeatletter
\newlabel{varieties-lemma-check-invertible-sheaf-trivial}{{0B40}{source}{}{sourcepointer-5}{}}
\newlabel{varieties-lemma-check-invertible-sheaf-trivial@cref}{{[originallemma][0][]0B40}{[1][0][]source}{}{}{}}
\makeatother
\crefname{originallemma}{original lemma}{original lemmas}
\Crefname{originallemma}{Original lemma}{Original lemmas}
\makeatletter
\newlabel{varieties-lemma-degree-pullback-map-proper-curves}{{0AYZ}{source}{}{sourcepointer-6}{}}
\newlabel{varieties-lemma-degree-pullback-map-proper-curves@cref}{{[originallemma][0][]0AYZ}{[1][0][]source}{}{}{}}
\makeatother
\crefname{originallemma}{original lemma}{original lemmas}
\Crefname{originallemma}{Original lemma}{Original lemmas}
\makeatletter
\newlabel{varieties-lemma-inseparable-deg-p-smooth}{{0CCF}{source}{}{sourcepointer-7}{}}
\newlabel{varieties-lemma-inseparable-deg-p-smooth@cref}{{[originallemma][0][]0CCF}{[1][0][]source}{}{}{}}
\makeatother
\title{Dichotomy for linear systems in characteristic p: generically etale pencil or small-degree system on the Frobenius twist}
\author{The Stacks Project Authors}
\date{stacks-project:curves.tex:local-snapshot-20261012; GNU FDL 1.2 or later}
\begin{document}
\maketitle
\noindent\textit{Editorial note.} All delivered mathematical text is the author's own, copied from the pinned Stacks Project chapter file curves.tex; no statement or proof step was rewritten, reordered or removed.\par\medskip
\noindent\textit{Editorial note.} The counted proof is the author's complete argument: the pencil test, the all-quotients factorization with the normal-curve extension, the pullback N with its independent sections, and the full degree computation are all retained verbatim.\par\medskip
\noindent\textit{Editorial note.} The label lemma-inseparable-linear-system names the counted statement, and the proof environment immediately following it in the source is delivered verbatim as its proof.\par\medskip
\noindent\textit{Editorial note.} The counted statement is carried by the author's own lemma environment in the chapter source.\par\medskip
\noindent\textit{Allowed standard prerequisites (editorial).}\begin{itemize}
\item Morphisms defined by pencils and generic etaleness. Same chapter, Lemmas lemma-linear-series and lemma-generically-etale; Proposition proposition-unwind-morphism-smooth. Identifies the rational function s/t of a pencil with the induced map to P\textasciicircum{}1, and shows that failure of generic etaleness forces f into k(X\textasciicircum{}\{(p)\}).
\item Extending maps over normal curves. Same chapter, Lemma lemma-extend-over-normal-curve. Extends the factored map from the affine chart D\_+(T\_0) to all of the smooth curve X\textasciicircum{}\{(p)\}.
\item Degree bookkeeping under relative Frobenius. Varieties, Lemmas varieties-lemma-check-invertible-sheaf-trivial, varieties-lemma-degree-pullback-map-proper-curves and varieties-lemma-inseparable-deg-p-smooth. Justifies L' = F\textasciicircum{}*N, the triviality check on components, and deg(F) = p for the relative Frobenius of a smooth curve.
\end{itemize}\medskip
\ifdefined\hypertarget\hypertarget{sourcepointer-1}{}\fi
\noindent\textit{Source reference.} Original reference 0BXZ: sources/stacks-curves/pinned-original-source.tex, source lines 287--287. Prerequisite pointer invoked at proof line(s) 2949; it names a standard fact of the same chapter that the counted argument uses, and every step of the counted proof built on it is retained verbatim in the delivered body.\par\medskip
\ifdefined\hypertarget\hypertarget{sourcepointer-2}{}\fi
\noindent\textit{Source reference.} Original reference 0C1C: sources/stacks-curves/pinned-original-source.tex, source lines 2564--2564. Prerequisite pointer invoked at proof line(s) 2922; it names a standard fact of the same chapter that the counted argument uses, and every step of the counted proof built on it is retained verbatim in the delivered body.\par\medskip
\ifdefined\hypertarget\hypertarget{sourcepointer-3}{}\fi
\noindent\textit{Source reference.} Original reference 0CCP: sources/stacks-curves/pinned-original-source.tex, source lines 583--583. Prerequisite pointer invoked at proof line(s) 2921; it names a standard fact of the same chapter that the counted argument uses, and every step of the counted proof built on it is retained verbatim in the delivered body.\par\medskip
\ifdefined\hypertarget\hypertarget{sourcepointer-4}{}\fi
\noindent\textit{Source reference.} Original reference 0CD2: sources/stacks-curves/pinned-original-source.tex, source lines 2888--2888. Prerequisite pointer invoked at proof line(s) 2934; it names a standard fact of the same chapter that the counted argument uses, and every step of the counted proof built on it is retained verbatim in the delivered body.\par\medskip
\ifdefined\hypertarget\hypertarget{sourcepointer-5}{}\fi
\noindent\textit{Source reference.} Original lemma 0B40: sources/stacks-chow/tags.txt, source lines 13983--13983. Prerequisite pointer invoked at proof line(s) 2965; it names a standard fact of a sibling chapter that the counted argument uses, and every step of the counted proof built on it is retained verbatim in the delivered body.\par\medskip
\ifdefined\hypertarget\hypertarget{sourcepointer-6}{}\fi
\noindent\textit{Source reference.} Original lemma 0AYZ: sources/stacks-chow/tags.txt, source lines 13807--13807. Prerequisite pointer invoked at proof line(s) 2966; it names a standard fact of a sibling chapter that the counted argument uses, and every step of the counted proof built on it is retained verbatim in the delivered body.\par\medskip
\ifdefined\hypertarget\hypertarget{sourcepointer-7}{}\fi
\noindent\textit{Source reference.} Original lemma 0CCF: sources/stacks-chow/tags.txt, source lines 15505--15505. Prerequisite pointer invoked at proof line(s) 2967; it names a standard fact of a sibling chapter that the counted argument uses, and every step of the counted proof built on it is retained verbatim in the delivered body.\par\medskip
\begin{lemma}
\label{lemma-inseparable-linear-system}
Let $k$ be a field of characteristic $p > 0$. Let $X$ be a smooth proper
curve over $k$. Let $(\mathcal{L}, V)$ be a $\mathfrak g^r_d$ with $r \geq 1$.
Then one of the following two is true
\begin{enumerate}
\item there exists a $\mathfrak g^1_d$ whose corresponding morphism
$X \to \mathbf{P}^1_k$ (Lemma \ref{lemma-linear-series})
is generically \'etale (i.e., is as in Lemma \ref{lemma-generically-etale}), or
\item there exists a $\mathfrak g^r_{d'}$ on $X^{(p)}$ where
$d' \leq d/p$.
\end{enumerate}
\end{lemma}

\begin{proof}
Pick two $k$-linearly independent elements $s, t \in V$.
Then $f = s/t$ is the rational function defining the morphism
$X \to \mathbf{P}^1_k$ corresponding to the linear series
$(\mathcal{L}, ks + kt)$. If this morphism is not
generically \'etale, then $f \in k(X^{(p)})$ by
Proposition \ref{proposition-unwind-morphism-smooth}.
Now choose a basis $s_0, \ldots, s_r$ of $V$ and let
$\mathcal{L}' \subset \mathcal{L}$ be the invertible sheaf
generated by $s_0, \ldots, s_r$. Set $f_i = s_i/s_0$ in $k(X)$.
If for each pair $(s_0, s_i)$ we have $f_i \in k(X^{(p)})$, then
the morphism
$$
\varphi = \varphi_{(\mathcal{L}', (s_0, \ldots, s_r)} :
X
\longrightarrow
\mathbf{P}^r_k = \text{Proj}(k[T_0, \ldots, T_r])
$$
factors through $X^{(p)}$ as this is true over the affine open
$D_+(T_0)$ and we can extend the morphism over the affine part
to the whole of the smooth curve $X^{(p)}$ by
Lemma \ref{lemma-extend-over-normal-curve}.
Introducing notation, say we have the factorization
$$
X \xrightarrow{F_{X/k}} X^{(p)} \xrightarrow{\psi} \mathbf{P}^r_k
$$
of $\varphi$. Then $\mathcal{N} = \psi^*\mathcal{O}_{\mathbf{P}^1_k}(1)$
is an invertible $\mathcal{O}_{X^{(p)}}$-module with
$\mathcal{L}' = F_{X/k}^*\mathcal{N}$ and with
$\psi^*T_0, \ldots, \psi^*T_r$ $k$-linearly independent
(as they pullback to $s_0, \ldots, s_r$ on $X$).
Finally, we have
$$
d = \deg(\mathcal{L}) \geq \deg(\mathcal{L}') =
\deg(F_{X/k}) \deg(\mathcal{N}) = p \deg(\mathcal{N})
$$
as desired. Here we used Varieties, Lemmas
\ref{varieties-lemma-check-invertible-sheaf-trivial},
\ref{varieties-lemma-degree-pullback-map-proper-curves}, and
\ref{varieties-lemma-inseparable-deg-p-smooth}.
\end{proof}

\end{document}
